\section{Sentiment Reversal}
\label{section:sentiment}




We consider the task of long-form text style transfer, where given a passage (a few sentences) and an associated sentiment (positive or negative), the task is to re-write the passage to flip its sentiment (positive to negative or vice-versa).
While a large body of work on style transfer is directed at sentence-level sentiment transfer~\citep{li-etal-2018-delete,prabhumoye-etal-2018-style}, we focus on transferring the sentiment of entire reviews, making the task challenging and providing opportunities for iterative improvements.


\paragraph{Instantiating \ours for sentiment reversal}
We instantiate \ours for this task following the high-level description of the framework shared in \Cref{sec:method}.
Recall that our requires three components: \initmod to generate an initial output, \fbmod to generate feedback on the initial output, and \itermod for improving the output based on the feedback.


\ours is implemented in a complete few-shot setup, where each module (\init, \fb, \iter) is implemented as few-shot prompts.
We execute the self-improvement loop for a maximum of $k=4$ iterations.
The iterations continue until the target sentiment is reached.



\subsection{Details}


\paragraph{Evaluation} 
Given an input and a desired sentiment level, we generate outputs \ours and the baselines.
Then, we measure the \% of times output from each setup was preferred to better align with the desired sentiment level~(see \Cref{sec:method} for more details).



We also experiment with standard text-classification metric. 
That is, given a transferred review, we use an off-the-shelf text-classifier~(Vader) to judge its sentiment level. 
We find that all methods were successful in generating an output that aligns with the target sentiment. 
For instance, when the target sentiment was positive, both \direct with \gptlatest and \ours generates sentences that have a positive sentiment (100\% classification accuracy). 
With the negative target sentiment, the classification scores were 92\% for \direct and 93.6\% for \ours. 

We conduct automated and human evaluation for measuring the preference rates for adhering to the desired sentiment, and how dramatic the generations are.
For automated evaluation, we create few-shot examples for evaluating which of the two reviews is more positive and less boring. We use a separate prompt for each task.
The examples are depicted in \Cref{fig:prompt:sentiment:init} for initialization, \Cref{fig:prompt:sentiment:fb} for feedback generation, and \Cref{fig:prompt:sentiment:refine} for refinement.
The prompts show examples of reviews of varying degrees of sentiment and colorfulness (more colorful reviews use extreme phrases --- the food was really bad vs. I wouldn't eat it if they pay me.).
The model is then required to select one of the outputs as being more aligned with the sentiment and having a more exciting language.
We report the preference rates: the \% of times a variant was preferred by the model over the outputs generated by \ours.



\paragraph{Pin-pointed feedback}

A key contribution of our method is supplying chain-of-thought prompting style feedback.
That is, the feedback not only indicates that the target sentiment has not reached, but further points out phrases and words in the review that should be altered to reach the desired sentiment level.
We experiment with an ablation of our setup where the feedback module simply says ``something is wrong.''
In such cases, for sentiment evaluation, the output from \ours were preferred 73\% of the time (down from 85\% with informative feedback).
For dramatic response evaluation, we found that the preference rate went down drastically to 58.92\%, from 80.09\%.
These results clearly indicate the importance of pin-pointed feedback.







\paragraph{Evaluation}

We evaluate the task using \gptf. Specifically, we use the following prompt:





When both win, we add winning rate to either. 